\section*{Chapter \ref{ch:g12:lists-stacks-queues} --- Linked Lists, Stacks and Queues}

\begin{solution}{exo:g12:lists-stacks-queues:1}
The \omterm{def:g12:lists-stacks-queues:stack}{pops} return $3$, then $4$, then $2$; the \omterm{def:g12:lists-stacks-queues:stack}{stack} holds $1$ alone.
\end{solution}

\begin{solution}{exo:g12:lists-stacks-queues:2}
The \omterm{def:g12:lists-stacks-queues:queue}{dequeues} return \code{a}, then \code{b}; the \omterm{def:g12:lists-stacks-queues:queue}{queue} holds \code{c} then \code{d}.
\end{solution}

\begin{solution}{exo:g12:lists-stacks-queues:3}
Bottom first: \texttt{\{}, \texttt{\{[}, \texttt{\{[(}, then \texttt{\{[} after \texttt{)}, \texttt{\{} after
\texttt{]}, empty after \texttt{\}}, and \texttt{(} after the last character. The string ends with a bracket still open: \code{False}.
\end{solution}

\begin{solution}{exo:g12:lists-stacks-queues:4}
Bottom first: $[5]$, $[5, 1]$, $[5, 1, 2]$, then \code{+} gives $[5, 3]$, $[5, 3, 4]$, \code{*} gives $[5, 12]$,
\code{+} gives $[17]$, $[17, 3]$, and \code{-} gives $[14]$. The value is $14$, that is $5 + (1 + 2) \times 4 - 3$.
\end{solution}

\begin{solution}{exo:g12:lists-stacks-queues:5}
In the list \omterm{def:g12:lists-stacks-queues:stack}{stack}, \code{return len(self._items)}. The linked \omterm{def:g12:lists-stacks-queues:stack}{stack} keeps a counter: the \omterm{def:g12:classes-and-objects:constructor}{constructor} sets
\code{self._size = 0}, \code{push} adds $1$, \code{pop} subtracts $1$ after a successful \omterm{def:g12:lists-stacks-queues:stack}{pop}, and \code{size} returns
it. The count is then a second invariant of the class: \code{_size} is the number of \omterm{def:g12:lists-stacks-queues:linked-list}{cells}.
\end{solution}

\begin{solution}{exo:g12:lists-stacks-queues:6}
The list \omterm{def:g12:lists-stacks-queues:stack}{stack} would raise \code{IndexError: pop from empty list}, from Python's list; the linked \omterm{def:g12:lists-stacks-queues:stack}{stack} would raise
\texttt{AttributeError: \textquotesingle{}NoneType\textquotesingle{} object has no attribute \textquotesingle{}value\textquotesingle{}}, since \code{_head} is \code{None}. Two
\omterm{def:g11:modules-and-interfaces:interface}{implementations} of one interface would fail in two different ways, and the second message speaks of the
\omterm{def:g11:modules-and-interfaces:interface}{implementation}, not of the \omterm{def:g12:lists-stacks-queues:stack}{stack}. The test makes both raise the error the interface promises, with a message a user
understands.
\end{solution}

\begin{solution}{exo:g12:lists-stacks-queues:7}
\smallskip\noindent\hspace*{2em}\code{def length_rec(cell):}

\noindent\hspace*{4em}\code{return 0 if cell is None else 1 + length_rec(cell.next)}

\smallskip\noindent\hspace*{2em}\code{def length_iter(cell):}

\noindent\hspace*{4em}\code{n = 0}

\noindent\hspace*{4em}\code{while cell is not None:}

\noindent\hspace*{6em}\code{n, cell = n + 1, cell.next}

\noindent\hspace*{4em}\code{return n}

\smallskip\noindent The recursive one needs one frame per cell, and fails beyond Python's recursion limit
(chapter \ref{ch:g12:recursion}); the loop does not.
\end{solution}

\begin{solution}{exo:g12:lists-stacks-queues:8}
After the \omterm{def:g12:lists-stacks-queues:queue}{dequeue}, \code{_head} is \code{None} but \code{_tail} still points at the cell of $1$. Enqueueing $2$ finds
\code{_tail} not \code{None}, so it links the new cell after the cell of $1$, which is no longer in the \omterm{def:g12:lists-stacks-queues:queue}{queue}, and
leaves \code{_head} at \code{None}. The $2$ is lost: nothing in the \omterm{def:g12:lists-stacks-queues:queue}{queue} leads to it. The next \omterm{def:g12:lists-stacks-queues:queue}{dequeue} reads
\code{self._head.value} and fails with \code{AttributeError}. The repair: in \code{dequeue}, set \code{_tail} to
\code{None} when \code{_head} becomes \code{None}.
\end{solution}

\begin{solution}{exo:g12:lists-stacks-queues:9}
\omterm{def:g12:lists-stacks-queues:queue}{Enqueues}: $1$, $2$, $3$ moves. The first \omterm{def:g12:lists-stacks-queues:queue}{dequeue} finds the \omterm{def:g12:lists-stacks-queues:stack}{out-stack} empty, transfers $3$ elements ($6$) and \omterm{def:g12:lists-stacks-queues:stack}{pops}
$1$ ($7$). \omterm{def:g12:lists-stacks-queues:queue}{Enqueue} $4$: $8$. \omterm{def:g12:lists-stacks-queues:queue}{Dequeue} $2$: $9$; \omterm{def:g12:lists-stacks-queues:queue}{dequeue} $3$: $10$. The last \omterm{def:g12:lists-stacks-queues:queue}{dequeue} finds the \omterm{def:g12:lists-stacks-queues:stack}{out-stack} empty again,
transfers $4$ ($11$) and \omterm{def:g12:lists-stacks-queues:stack}{pops} it ($12$). Four elements, $12 = 3 \times 4$ moves.
\end{solution}

\begin{solution}{exo:g12:lists-stacks-queues:10}
\smallskip\noindent\hspace*{2em}\code{def reverse(head):}

\noindent\hspace*{4em}\code{previous = None}

\noindent\hspace*{4em}\code{while head is not None:}

\noindent\hspace*{6em}\code{following = head.next}

\noindent\hspace*{6em}\code{head.next = previous}

\noindent\hspace*{6em}\code{previous, head = head, following}

\noindent\hspace*{4em}\code{return previous}

\smallskip\noindent Invariant of the loop: \code{previous} is the reversed list of the \omterm{def:g12:lists-stacks-queues:linked-list}{cells} already visited, and
\code{head} the rest, untouched. One pass, no new cell.
\end{solution}

\begin{solution}{exo:g12:lists-stacks-queues:11}
\code{2 3 + 7 4 - * 5 /}. The \omterm{def:g12:lists-stacks-queues:stack}{stack}, bottom first: $[2]$, $[2, 3]$, $[5]$, $[5, 7]$, $[5, 7, 4]$, $[5, 3]$, $[15]$,
$[15, 5]$, $[3]$: the value is $3$.
\end{solution}

\begin{solution}{exo:g12:lists-stacks-queues:12}
Follow one element. It is pushed on the \omterm{def:g12:lists-stacks-queues:stack}{in-stack} once, by its \omterm{def:g12:lists-stacks-queues:queue}{enqueue}: one move. Only a transfer takes elements off
the \omterm{def:g12:lists-stacks-queues:stack}{in-stack}, and a transfer puts each one on the \omterm{def:g12:lists-stacks-queues:stack}{out-stack}: one move. Nothing ever puts an element back on the
\omterm{def:g12:lists-stacks-queues:stack}{in-stack}, so an element is transferred at most once. Only a \omterm{def:g12:lists-stacks-queues:queue}{dequeue} takes an element off the \omterm{def:g12:lists-stacks-queues:stack}{out-stack}, and it
leaves the \omterm{def:g12:lists-stacks-queues:queue}{queue}: one move. So each element makes at most three moves. Every move counted by the class is one of
these three, charged to one element; with $n$ elements enqueued, the total is at most $3n$, whatever the order of
the operations. A single \omterm{def:g12:lists-stacks-queues:queue}{dequeue} may make many moves, but the total stays in proportion to $n$.
\end{solution}

\begin{solution}{exo:g12:lists-stacks-queues:13}
\smallskip\noindent\hspace*{2em}\code{def enqueue(self, x):}

\noindent\hspace*{4em}\code{if self._count == len(self._slots):}

\noindent\hspace*{6em}\code{raise OverflowError("the queue is full")}

\noindent\hspace*{4em}\texttt{self.\_slots[(self.\_first + self.\_count) \% len(self.\_slots)] = x}

\noindent\hspace*{4em}\code{self._count += 1}

\smallskip\noindent\hspace*{2em}\code{def dequeue(self):}

\noindent\hspace*{4em}\code{if self._count == 0:}

\noindent\hspace*{6em}\code{raise IndexError("dequeue from an empty queue")}

\noindent\hspace*{4em}\code{x = self._slots[self._first]}

\noindent\hspace*{4em}\texttt{self.\_first = (self.\_first + 1) \% len(self.\_slots)}

\noindent\hspace*{4em}\code{self._count -= 1}

\noindent\hspace*{4em}\code{return x}

\smallskip\noindent The \omterm{def:g12:classes-and-objects:constructor}{constructor} sets \code{_slots = [None] * capacity}, \code{_first = 0} and \code{_count = 0}.
The remainder by the capacity makes the index wrap from the last slot to slot $0$; no element ever moves.
\end{solution}

\begin{solution}{exo:g12:lists-stacks-queues:14}
\smallskip\noindent\hspace*{2em}\code{PRIORITY = {"+": 1, "-": 1, "*": 2, "/": 2}}

\smallskip\noindent\hspace*{2em}\code{def to_postfix(expression):}

\noindent\hspace*{4em}\code{out, ops = [], Stack()}

\noindent\hspace*{4em}\code{for token in expression.split():}

\noindent\hspace*{6em}\code{if token in PRIORITY:}

\noindent\hspace*{8em}\code{while (not ops.is_empty() and ops.top() != "("}

\noindent\hspace*{12em}\code{and PRIORITY[ops.top()] >= PRIORITY[token]):}

\noindent\hspace*{10em}\code{out.append(ops.pop())}

\noindent\hspace*{8em}\code{ops.push(token)}

\noindent\hspace*{6em}\code{elif token == "(":}

\noindent\hspace*{8em}\code{ops.push(token)}

\noindent\hspace*{6em}\code{elif token == ")":}

\noindent\hspace*{8em}\code{while ops.top() != "(":}

\noindent\hspace*{10em}\code{out.append(ops.pop())}

\noindent\hspace*{8em}\code{ops.pop()}

\noindent\hspace*{6em}\code{else:}

\noindent\hspace*{8em}\code{out.append(token)}

\noindent\hspace*{4em}\code{while not ops.is_empty():}

\noindent\hspace*{6em}\code{out.append(ops.pop())}

\noindent\hspace*{4em}\code{return " ".join(out)}

\smallskip\noindent ``Greater than or equal'' makes operators of equal priority leave from the left:
\code{10 - 2 - 3} gives \code{10 2 - 3 -}, which is $5$, and not $10 - (2 - 3)$. The tests check the expression of
exercise \ref{exo:g12:lists-stacks-queues:11}, and that \code{1 + 2 * 3 - 8 / 4} evaluates to $5$.
\end{solution}

\begin{solution}{exo:g12:lists-stacks-queues:15}
Only the bracket characters count here; other characters do nothing to the \omterm{def:g12:lists-stacks-queues:stack}{stack}.

\emph{Balanced strings are accepted.} We prove more, by induction on the length: started with any \omterm{def:g12:lists-stacks-queues:stack}{stack} $S$, the loop
reads a balanced string $w$ without returning \code{False} and ends with the \omterm{def:g12:lists-stacks-queues:stack}{stack} $S$ again. The empty string does
nothing. If $w = (u)v$, the \code{(} is pushed; $u$, shorter, leaves the \omterm{def:g12:lists-stacks-queues:stack}{stack} as $S$ with \code{(} on top; the
\code{)} matches that top and \omterm{def:g12:lists-stacks-queues:stack}{pops} it, giving $S$; $v$, shorter, leaves $S$. The same for $[u]v$. Started with an
empty \omterm{def:g12:lists-stacks-queues:stack}{stack}, the function therefore ends with an empty \omterm{def:g12:lists-stacks-queues:stack}{stack} and returns \code{True}.

\emph{Accepted strings are balanced}, by induction on the length. The empty string is balanced. Let $w$ be accepted
and not empty. Its first character is an opening bracket, since a closing one would meet an empty \omterm{def:g12:lists-stacks-queues:stack}{stack}. That bracket
lies at the bottom of the \omterm{def:g12:lists-stacks-queues:stack}{stack}, and is popped at some point, since the \omterm{def:g12:lists-stacks-queues:stack}{stack} ends empty: by a closing bracket of the
same kind, at some position. Call $u$ the characters strictly between, and $v$ those after. While $u$ is read, the
bottom bracket is never reached, so $u$ is read exactly as from an empty \omterm{def:g12:lists-stacks-queues:stack}{stack}, and ends with the \omterm{def:g12:lists-stacks-queues:stack}{stack} holding only
that bracket: $u$ is accepted, and balanced by induction. After the matching closing bracket the \omterm{def:g12:lists-stacks-queues:stack}{stack} is empty, and
$v$ is read from an empty \omterm{def:g12:lists-stacks-queues:stack}{stack} and accepted: balanced by induction. So $w = (u)v$ or $[u]v$ is balanced.
\end{solution}

\begin{solution}{pb:g12:lists-stacks-queues:1}
\textbf{1.} The \omterm{def:g12:classes-and-objects:constructor}{constructor}, \code{enqueue}, \code{dequeue} and \code{is_empty}. It also reads \code{moves}, a counter
added for this chapter, which is not part of the \omterm{def:g12:lists-stacks-queues:queue}{queue}'s interface.

\textbf{2.} A class is a value: \code{simulate} receives it and calls it to build a \omterm{def:g12:lists-stacks-queues:queue}{queue}. It then uses only the
interface, which both classes offer with the same behaviour.

\textbf{3.} The minute of each arrival. The only thing the program needs to know about a customer is how long they
waited, and that is the minute of service minus the minute of arrival.

\textbf{4.} The \omterm{def:g12:lists-stacks-queues:adt}{abstract data type} has no operation giving the number of elements. Counting arrivals and departures
answers in one step, where walking the \omterm{def:g12:lists-stacks-queues:queue}{queue} would cost a step per person.

\textbf{5.} On an empty \omterm{def:g12:lists-stacks-queues:queue}{queue}, no link between \omterm{def:g12:lists-stacks-queues:linked-list}{cells}: \code{_head} and \code{_tail} are both set to the new cell. On
a \omterm{def:g12:lists-stacks-queues:queue}{queue} of three, the link of the old last cell, and \code{_tail}.

\textbf{6.} With this \code{enqueue}, which tests emptiness on \code{_head}, the \omterm{def:g12:lists-stacks-queues:queue}{queue} would still work. But
\code{_tail} would keep pointing at a cell that has left the \omterm{def:g12:lists-stacks-queues:queue}{queue}, so the invariant ``\code{_tail} is \code{None}
exactly when the \omterm{def:g12:lists-stacks-queues:queue}{queue} is empty'' would be false, and any method relying on it would fail, as in exercise
\ref{exo:g12:lists-stacks-queues:8}.

\textbf{7.} When the \omterm{def:g12:lists-stacks-queues:stack}{out-stack} is empty. It transfers every element of the \omterm{def:g12:lists-stacks-queues:stack}{in-stack}.

\textbf{8.} \omterm{def:g12:lists-stacks-queues:queue}{Enqueue} $1$, $2$, $3$, then \omterm{def:g12:lists-stacks-queues:queue}{dequeue} three times. The first \omterm{def:g12:lists-stacks-queues:queue}{dequeue} makes $4$ moves: $3$ transfers and one
\omterm{def:g12:lists-stacks-queues:stack}{pop}.

\textbf{9.} For both, the linked \omterm{def:g12:lists-stacks-queues:queue}{queue}: it makes fewer moves, and every operation takes a time that does not depend
on the length, where a \omterm{def:g12:lists-stacks-queues:queue}{dequeue} of the \omterm{def:g12:lists-stacks-queues:stack}{two-stack} \omterm{def:g12:lists-stacks-queues:queue}{queue} can transfer the whole line at once.

\textbf{10.} So that the run is the same day every time: the output printed in the book is reproduced exactly, and
both \omterm{def:g11:modules-and-interfaces:interface}{implementations} are compared on the same arrivals.

\textbf{11.} The \omterm{def:g12:lists-stacks-queues:queue}{queue} draws no random numbers, so the generator gives the same numbers in the same order to both
runs; and both are \omterm{def:g12:lists-stacks-queues:queue}{queues}, giving back the same arrival minutes in the same order. Every decision of the program is
then the same.

\textbf{12.} The clerk is already busy most of the time. If services were longer than the time between arrivals,
people would arrive faster than they leave, and the line would grow all day.

\textbf{13.} The doors close at minute $480$, but the people inside are served. The loop stops at the first minute
from $480$ on at which the line is empty.

\textbf{14.} The \omterm{def:g12:lists-stacks-queues:queue}{queue} holds each customer's arrival minute; when one is served, at minute \code{minute}, the wait is
\code{minute - arrived}, and the longest is kept.

\textbf{15.} $114$ customers; at most $5$ people waiting at once; the longest wait $13$ minutes.

\textbf{16.} $2$: one \omterm{def:g12:lists-stacks-queues:queue}{enqueue} and one \omterm{def:g12:lists-stacks-queues:queue}{dequeue} per customer, and all $114$ were served.

\textbf{17.} $3$: every customer is enqueued, and, being served, has been transferred once and popped once. The bound
of exercise \ref{exo:g12:lists-stacks-queues:12} is reached because nobody is left in line.

\textbf{18.} A move of either \omterm{def:g12:lists-stacks-queues:queue}{queue} creates or unlinks one cell, a few steps of the same kind. The \omterm{def:g12:lists-stacks-queues:stack}{two-stack} \omterm{def:g12:lists-stacks-queues:queue}{queue} does
$342$ moves against $228$: half as much work again, but the same order of cost, in proportion to the number of
customers.

\textbf{19.} The people still in line would not be dequeued: fewer than $2$ moves per customer for the linked \omterm{def:g12:lists-stacks-queues:queue}{queue},
fewer than $3$ for the \omterm{def:g12:lists-stacks-queues:stack}{two-stack} \omterm{def:g12:lists-stacks-queues:queue}{queue}, where those still in the \omterm{def:g12:lists-stacks-queues:stack}{in-stack} would have made only one.

\textbf{20.} \textbf{Over the day, $114$ customers, a longest line of $5$ and a longest wait of $13$ minutes, the same
with both \omterm{def:g12:lists-stacks-queues:queue}{queues}; the linked \omterm{def:g12:lists-stacks-queues:queue}{queue} makes $228$ moves, $2$ per customer, the \omterm{def:g12:lists-stacks-queues:stack}{two-stack} \omterm{def:g12:lists-stacks-queues:queue}{queue} $342$, exactly $3$ per
customer.}
\end{solution}
